Berdasarkan latar belakang masalah yang telah diuraikan, rumusan masalah dalam penelitian ini adalah sebagai berikut:
\begin{enumerate}
    \item Bagaimana merancang alur kerja (\textit{workflow}) pemindahan data historis mahasiswa secara otomatis dari basis data produksi ke basis data arsip, guna mengoptimalkan kinerja eksekusi kueri dan efisiensi ruang penyimpanan pada Pusdatin UNTAR?

    \item Bagaimana membangun portal \textit{web} berbasis ASP.NET yang memfasilitasi pencarian dan penarikan data secara \textit{on-demand} (\textit{on-demand retrieval}) secara fleksibel lintas basis data, baik basis data aktif maupun basis data arsip, tanpa mengganggu kinerja sistem produksi?

    \item Bagaimana merancang dan mengimplementasikan arsitektur \textit{Client-Worker} terdistribusi berbasis HTTP API JSON (\textit{agent.ashx}) yang dieksekusi melalui SQL Server Agent dan skrip PowerShell pada mesin eksternal terpisah, guna menjalankan antrean tugas (\textit{task queue}) pengarsipan data tanpa membebani infrastruktur server utama kampus?
\end{enumerate}
